\documentclass[12pt]{article}
\usepackage[utf8]{inputenc}
\usepackage[T1]{fontenc}
\usepackage[margin=1in]{geometry}
\usepackage{setspace}
\usepackage{xcolor}
\usepackage[colorlinks=true, urlcolor=blue!60!black, linkcolor=blue!60!black]{hyperref}
\setlength{\parindent}{0em}
\setlength{\parskip}{0.6em}
\renewcommand{\baselinestretch}{1.0}
\begin{document}
\begin{center}
{\LARGE \textit{Allen v.\ Cooper}, 589 U.S.\ 248 (2020)\par}
\vspace{0.5em}
{\large Case Brief --- Two-Page Edition\par}
\end{center}

\section{Memory Jogger}
North Carolina posted a videographer's footage of Blackbeard's shipwreck online; the Court unanimously held that Congress could not abrogate state immunity from copyright damages under Article~I or Section~5 (\href{https://supreme.justia.com/cases/federal/us/589/18-877/}{Justia}).

\section{Detailed Case Facts}
Frederick Allen documented the salvage of the \textit{Queen Anne's Revenge}, which belonged to North Carolina, and registered copyrights in his videos and photographs. After a \$15,000 settlement, the State again posted five videos and published a photograph, and enacted ``Blackbeard's Law'' making such materials ``public record.'' Allen sued for damages under the Copyright Remedy Clarification Act (CRCA), which strips state immunity.

\section{Procedural History}
The district court upheld the CRCA under Section~5; the Fourth Circuit reversed under \textit{Florida Prepaid}; the Supreme Court affirmed.

\section{Judicial Votes}
Affirmed, 9--0. Kagan wrote for the Court (Thomas joining all but two paragraphs); Thomas concurred in part and in the judgment; Breyer, joined by Ginsburg, concurred in the judgment.

\section{Holding}
The CRCA's abrogation is invalid. Article~I's Intellectual Property Clause cannot abrogate immunity (\textit{Seminole Tribe}, \textit{Florida Prepaid}), and \textit{Katz} is limited to bankruptcy. Under Section~5, blanket abrogation fails \textit{Boerne}'s ``congruence and proportionality'' test given a thin record of mostly negligent infringement.

\section{Analysis of Opinions}
\textbf{Majority:} ``\textit{Florida Prepaid} all but prewrote our decision today.'' Overruling it would require a ``special justification.'' Infringement violates due process only if intentional or reckless and unremedied by the State; the Oman Report found twelve possible instances, two apparently intentional. Congress may try a tailored statute. \textbf{Thomas:} rejects the ``special justification'' standard and asks whether copyrights are due process ``property.'' \textbf{Breyer:} the Court ``went astray'' in \textit{Seminole Tribe} but precedent controls.

\section{Examples: Future Applications}
\textbf{Same-Side:} a publisher's CRCA damages suit against a state university for a professor's scanned chapters is barred. \textbf{Opposite-Side:} an \textit{Ex parte Young} suit for prospective relief against a state official, or a damages suit against a city agency, may proceed. \textbf{Fence-Sitter:} a narrow statute reaching only willful state infringement without adequate state remedies, backed by a new record.

\section{Critique}
Vikram David Amar argues the Court measures the frequency of state infringement rather than the proportion of covered conduct that is unconstitutional, and faults the ``special justification'' rule. The ``good-for-one-clause-only'' reading of \textit{Katz} collapsed two years later in \textit{Torres}, and progressives fault the Court for creating a right without a damages remedy.

\section{Key Quotations}
\begin{quote}
``But everything in \textit{Katz} is about and limited to the Bankruptcy Clause; the opinion reflects what might be called bankruptcy exceptionalism.''
\end{quote}
\noindent --- Justice Kagan, Opinion of the Court.

\end{document}
